% !TEX root = ../main.tex
%% Devoir surveillé n° 3 — chapitres 7 à 10
%% Résultats vérifiés par verifications/eval_ds.py.

\sujetEval{DS3}{MPSI · Devoir surveillé}{Devoir surveillé \numero3}
  {Durée : \textbf{4 heures}\\Calculatrice interdite\\Chapitres 7 à 10}{4}
  {Exercice 1/Arithmétique/5 pts,Exercice 2/Groupes/5 pts,%
   Exercice 3/Calcul matriciel/5 pts,Exercice 4/Polynômes et fractions/5 pts}
\consigne{Les quatre exercices sont indépendants et peuvent être traités
dans l'ordre de son choix.}

\begin{exoEval}{5 points}{Arithmétique}
\begin{questions}
  \item Résoudre dans $\Z^{2}$ l'équation $21x+15y=9$.
  \item Déterminer les deux derniers chiffres de l'écriture décimale de
    $7^{2026}$.
  \item Montrer que pour tout $n\in\Z$, $n^{5}-n$ est divisible par $30$.
  \item Pour $n\in\N$, montrer que $(n^{2}+1)\wedge(n+1)\in\ens{1,2}$, et
    préciser quand il vaut $2$.
\end{questions}
\end{exoEval}

\begin{exoEval}{5 points}{Groupes}
Soit $n\in\N^{*}$ et $\omega=\mathrm{e}^{2\mathrm{i}\pi/n}$.
\begin{questions}
  \item Montrer que $\mathbb{U}_n$ est un sous-groupe de $(\Cx^{*},\times)$,
    puis que $\varphi:k\mapsto\omega^{k}$ est un morphisme surjectif de
    $(\Z,+)$ sur $\mathbb{U}_n$ dont on déterminera le noyau.
  \item Montrer que les sous-groupes de $(\Z,+)$ sont les $m\Z$,
    $m\in\N$. \emph{Considérer le plus petit élément strictement positif.}
  \item Montrer que l'ordre de $\omega^{k}$ est $\dfrac{n}{n\wedge k}$.
    Quels éléments de $\mathbb{U}_{12}$ engendrent ce groupe ?
\end{questions}
\end{exoEval}

\begin{exoEval}{5 points}{Calcul matriciel}
\begin{questions}
  \item Soit $A=\begin{pmatrix}1&1&1\\0&1&1\\0&0&1\end{pmatrix}$. En
    écrivant $A=I_3+N$, calculer $A^{n}$ pour tout $n\in\N$.
  \item Montrer que $A$ est inversible et que la formule précédente reste
    vraie pour $n=-1$.
  \item Résoudre, selon la valeur du réel $m$, le système
    \[
    \begin{cases}x+y+mz=1\\ x+my+z=1\\ mx+y+z=1\end{cases}
    \]
\end{questions}
\end{exoEval}

\begin{exoEval}{5 points}{Polynômes et fractions}
On note $P=X^{4}+2X^{3}+3X^{2}+2X+1$ et $F=\dfrac{1}{X^{3}-X}$.
\begin{questions}
  \item Développer $(X^{2}+X+1)^{2}$, puis factoriser $P$ dans $\Cx[X]$ et
    dans $\R[X]$. Retrouver la multiplicité de la racine $j$ à l'aide des
    dérivées.
  \item Pour $n\geqslant2$, déterminer le reste de la division euclidienne
    de $X^{n}$ par $(X-1)^{2}$.
  \item Décomposer $F$ en éléments simples, et en déduire une primitive de
    $x\mapsto F(x)$ sur $\intOO{1}{+\infty}$.
\end{questions}
\end{exoEval}
